\chapter{Getting Started}
\label{ch:getting-started}

There are three ways to use \jns{}, in increasing order of commitment:
calibrating the uncertainty threshold (no external program), evaluating
structures with the safe calculator (needs the models, falls back to ORCA
only when triggered), and running the models inside the SCINE engine (needs
the full stack).  Each path below is self-contained.

\section{Path 1: calibrating an uncertainty threshold}

The threshold that decides when the machine-learning potential is trusted is
a measured quantity, not a default.  Calibrating it needs two arrays
collected on a reference set: the committee disagreement \code{uq} of each
configuration, and the true error \code{err} of that configuration against a
reference calculation (for the force criterion, the per-atom force error in
eV/\AA).  The calibration itself is pure \file{numpy} and runs without any
external program:

\begin{codeblock}
import numpy as np
from jnuscout.uq.calibration import calibrate_threshold, coverage_stats, spearman_rho

uq  = np.array([0.02, 0.05, 0.11, 0.30, 0.08, 0.45, 0.15, 0.60])
err = np.array([0.01, 0.03, 0.09, 0.40, 0.02, 0.55, 0.12, 0.70])  # eV/A

theta = calibrate_threshold(uq, err, epsilon_crit=0.1)
print(theta)                             # 0.1725 with these numbers
print(spearman_rho(uq, err))             # 0.98 (rank correlation, want > 0.6)
print(coverage_stats(uq, err, theta, epsilon_crit=0.1))
\end{codeblock}

With these numbers the script prints:

\begin{transcript}
0.1725
0.9761904761904762
{'n': 8, 'coverage': 0.625, 'fallback_rate': 0.375, 'n_dangerous': 4,
 'false_negative_rate': 0.25, 'nn_mae': 0.054000000000000006,
 'all_mae': 0.24000000000000002, 'theta': 0.1725}
\end{transcript}

\begin{readbox}[Reading the calibration output]
\code{theta} is the fallback threshold: a configuration whose disagreement
reaches it is handed to the reference method.  \code{spearman_rho} measures
whether the disagreement ranks configurations the way the true error does; a
committee worth using reaches well above 0.6.  \code{coverage_stats}
reports the share of configurations carried by the potential, the
complementary fallback rate, and the false-negative rate (dangerous
configurations that stayed below the threshold).  The example set is tiny ---
four dangerous configurations, one of them below \code{theta} --- so the
empirical rate is coarse here; the 5\,\% quantile construction is what bounds
it on a realistic reference set.  Chapter~8 gives the full protocol.
\end{readbox}

Two behaviours to keep in mind when wiring this into a script: thresholds
come in pairs with the accuracy target, so a stricter \code{epsilon_crit}
produces a lower \code{theta} and a higher fallback rate; and if the
reference set contains no dangerous configuration at all,
\code{calibrate_threshold} returns \code{inf} (never fall back), which the
caller should treat as an alarm rather than a result.

\section{Path 2: evaluating structures with the safe calculator}

This path needs the \code{[nnp]} extra.  ORCA is needed only if a fallback
is actually triggered:

\begin{codeblock}
from ase.io import read
from jnuscout.calculators import SafeNNPCalculator

atoms = read("molecule.xyz")
calc = SafeNNPCalculator(theta_force=0.30)        # theta from path 1
atoms.calc = calc
print(atoms.get_potential_energy(), atoms.get_forces())
print(calc.n_fallback, "/", calc.n_calls)
\end{codeblock}

Every call runs the committee once, forms the ensemble mean, and compares the
disagreement with the threshold.  Below the threshold the ensemble result is
returned; at or above it the structure is re-evaluated with ORCA L2 and that
result is returned instead.  The last line prints the number of fallbacks
over the total number of calculations, for example \code{0 / 1} for a single
well-behaved structure.  The ORCA fallback object is built lazily, so no ORCA
process exists unless a fallback happens.

\begin{notebox}[First run]
The default committee is MACE-MP-0, MACE-MPA-0 and AIMNet2-b973c, evaluated
on \code{device="cuda"}; pass \code{device="cpu"} on a machine without a GPU.
The first call loads the model weights, so it is much slower than the rest; on
an offline machine set \code{MACE_MP0_PATH} and \code{MACE_MPA0_PATH}
first (Chapter~\ref{ch:installation}).
\end{notebox}

\begin{notebox}[Caching]
Results follow the ASE calculator protocol: energy and forces of one
evaluation are cached together, and repeated requests on an unchanged
\code{Atoms} object reuse them without a new evaluation.  The counters
therefore count evaluations, not property accesses.
\end{notebox}

\section{Path 3: driving the SCINE engine}

To let the engine run the potential as a method family, the injection patch
must be active in every process of the job daemon before Puffin starts.  The
supported route is a \file{sitecustomize.py} on \code{PYTHONPATH}:

\begin{codeblock}
# sitecustomize.py, on PYTHONPATH
from jnuscout.scine import install_from_env
install_from_env()
\end{codeblock}

With \code{NNP_AFIR_PUFFIN=1} in the environment, the patch intercepts
calculator resolution for the declared method families and leaves every other
method family untouched, so DFTB3 and friends keep working; without the
variable the call is a no-op.  Start your Chemoton exploration as usual.  The
choke points, the environment variables and the bridge contract are the
subject of Chapter~7.

\section{A complete workflow}

The three paths compose into one loop.  A typical cycle:

\begin{enumerate}
  \item \emph{Pre-screen cheaply.}  L1 (GFN2-xTB) refines initial
  geometries and screens candidate structures at negligible cost.
  \item \emph{Explore with the potential.}  Chemoton drives the reaction
  network search while the committee evaluates structures; every geometry
  receives a disagreement score alongside its energy and forces.
  \item \emph{Route by uncertainty.}  Configurations whose disagreement
  reaches the calibrated threshold are re-evaluated at L2 (ORCA
  r2SCAN-D4/def2-SVP); their NNP answer is discarded.  Everything else stays
  on the fast path.
  \item \emph{Close the active-learning round.}  Sample frames along the
  explored paths, label them at L2, split by composition and group by
  reaction, fine-tune the potential, and verify on held-out reactions.
  \item \emph{Repeat.}  The next round starts from the fine-tuned potential,
  so the fallback rate at a fixed accuracy target falls as the potential
  learns the chemistry it is being asked about.
\end{enumerate}

Points that determine whether this loop helps or misleads: the training
composition sets the generality of the fine-tuned model (a single-composition
training set produces a specialist); reaction-type coverage is a separate
axis from composition coverage; and errors must be reported by system size,
not only in aggregate.  Chapter~9 documents the protocol and the pitfalls.

\section{Where to go next}

\begin{table}[htbp]
  \centering
  \small
  \begin{tabularx}{\linewidth}{@{}l X@{}}
    \toprule
    Chapter & Contents \\
    \midrule
    5 & Architecture: the three layers and what runs where \\
    \ref{ch:calculators} & Calculator reference: parameters, examples, caveats \\
    7 & SCINE integration: method-family injection and the bridge \\
    8 & Uncertainty: disagreement definitions and the calibration protocol \\
    9 & Active learning: training-set construction and validation \\
    10 & Scripts: command-line tools for calibration, labelling and database hygiene \\
    \ref{ch:configuration} & Example configuration files, field by field \\
    \bottomrule
  \end{tabularx}
  \caption{Reading paths from here.}
\end{table}
